% Results and implementation analysis owned by Ninh Ngoc Hung.
% The historical 60-second pilot is described in experiments.tex;
% this section reports the shared MIP formulation and the six-method run.

\subsection{MIP model size and scalability}
The CPLEX MIP and Gurobi MIP adapters use the same assignment-based
formulation. For a graph with $n$ vertices and $m$ undirected edges, the
implementation creates $n^2+1$ variables ($n^2$ binary assignment variables
and one integer objective variable) and $2n+mn^2$ constraints. The latter
count follows the implementation's edge--position generation loop, including
rows for position pairs whose cyclic distance is zero.

\begin{table}[t]
\centering
\small
\caption{Estimated size of the shared MIP formulation on the verified graph
set. Counts are computed before invoking either optimizer.}
\label{tab:mip-model-size}
\begin{tabular}{lrrrr}
\hline
Instance & $n$ & $m$ & Variables & Constraints \\
\hline
JGL009 & 9 & 32 & 82 & 2,610 \\
BCSSTK01 & 48 & 176 & 2,305 & 405,600 \\
BP 400 & 822 & 4,015 & 675,685 & 2,712,872,904 \\
\hline
\end{tabular}
\end{table}

The $n^2$ assignment variables and $mn^2$ edge--position constraints make
this direct formulation impractical on BP 400. The runner's one-million-
constraint preflight limit rejects that instance before allocating the MIP
model. This is a configured resource guard, not an optimizer result, timeout,
or infeasibility proof.

\subsection{Six-method benchmark outcomes}
The consolidated table contains 18 method--instance entries. Fourteen
have validated layouts, including the two CPLEX Studio reruns. The
remaining four were stopped by the BP 400 size guards.
Table~\ref{tab:six-method-rerun} uses a 300-second budget and one thread
per entry. Sixteen entries come from the original batch at commit
\texttt{9df7be6}. The JGL009 and BCSSTK01 entries for IBM CPLEX MIP
come from Studio 22.2 reruns with the same formulation and budget.

All six methods proved 4 on JGL009. On BCSSTK01, pairwise SAT and both CP
engines reached 12, order SAT reached 13, Gurobi MIP reached 15, and CPLEX
MIP reached 18. All six runs ended with a reported lower bound of 6 and
without an optimality proof. CPLEX reported a solver bound of 1, which
the independent degree bound strengthened to 6. On BP 400, pairwise SAT and
order SAT were rejected by their estimated-size guards, and both MIP
adapters were rejected by the one-million-constraint guard before solver
startup. CP-SAT and CP Optimizer returned validated incumbents of 399 and
263, respectively, with certified lower bound 147, but neither proved
optimality within the time limit.

Solver runtime is N/A for the four guarded BP~400 runs.

\subsection{Experimental limitations}
The historical main manifest contains three graphs, not the full 28-instance
candidate list in Table A.2 of Ren et al. The set provides a small dense,
medium sparse, and larger sparse example, but is too small for broad
statistical or scalability claims. BP 400 is an additional Harwell--Boeing /
SMTAPE instance, not one of the 28 Table A.2 candidates. The graph edge count
$m$ is the number of unique undirected edges after removing self-loops and
merging duplicate or opposite-direction entries; it is not the raw matrix
nonzero count.

Both MIP solvers produced validated layouts on four matrix graphs:
JGL009, BCSSTK01, IBM32 and BCSPWR01. Only JGL009 yielded an optimality
proof from both. The other three comparisons measure the layouts and
bounds reached within each budget. BP~400 provides a model-size estimate
only, and the separate BCSPWR02 experiment evaluated the SAT encodings.
